% © 2026 Ing. David Jaroš — CC BY-NC-ND 4.0
%
% This work is licensed under a Creative Commons Attribution-NonCommercial-NoDerivatives
% 4.0 International License (CC BY-NC-ND 4.0).
%
% License History: Earlier drafts (up to v0.3) were released under CC BY 4.0.
% From v0.4 onward, all material is released under CC BY-NC-ND 4.0 to protect
% the integrity of the theoretical work during ongoing academic development.
%
% See LICENSE.md for full license text.

% UBT-AI-PROVENANCE-BEGIN
% schema: ubt-ai-provenance/v1
% tier: B_machine_verified
% ai_assistance: disclosed
% human_review: machine-verification
% editorial_responsibility: Ing. David Jaroš
% policy: ../../AI_PROVENANCE.md
% notice: Machine-verified against named sources or verifiers; individual attestation is not claimed.
% UBT-AI-PROVENANCE-END

% -----------------------------------------------------------------------
% STATUS: HISTORICAL / COMPLEMENTARY — NOT THE CURRENT CANONICAL SU(3) DERIVATION
%
% This file contains the quantum-superposition/triqubit approach to deriving
% the SU(3) colour sector in UBT.  It is preserved for provenance and
% priority documentation.
%
% Key limitations (do not modify this notice):
%   1. The global group statement  U(3) ≅ U(1) × SU(3)  in Theorem 3.1 is
%      INCORRECT at the global level.  The correct statement is:
%          U(3) ≅ (SU(3) × U(1)) / Z_3
%      The Lie-algebra split  u(3) = su(3) ⊕ u(1)  is valid.
%   2. The unitary structure on the carrier was postulated, not derived.
%
% The current canonical derivation is:
%   canonical/su3_derivation/su3_stabilizer_exterior_fock.tex
% -----------------------------------------------------------------------

\ifdefined\INCLUDEMODE
\else
  \documentclass[12pt,a4paper]{article}
  \usepackage[a4paper, margin=2.5cm]{geometry}
  \usepackage{amsmath, amssymb, amsthm}
  \usepackage{mathtools}
  \usepackage{hyperref}
  \usepackage{xcolor}
  \usepackage{mdframed}
  \usepackage{booktabs}

  \newtheorem{theorem}{Theorem}[section]
  \newtheorem{lemma}[theorem]{Lemma}
  \newtheorem{proposition}[theorem]{Proposition}
  \newtheorem{corollary}[theorem]{Corollary}
  \theoremstyle{definition}
  \newtheorem{definition}[theorem]{Definition}
  \theoremstyle{remark}
  \newtheorem{remark}[theorem]{Remark}

  \newmdenv[backgroundcolor=yellow!20, linecolor=orange, linewidth=1.5pt]{gapbox}
  \newmdenv[backgroundcolor=green!15, linecolor=green!60!black, linewidth=1.5pt]{resultbox}
  \newmdenv[backgroundcolor=blue!8, linecolor=blue!50, linewidth=1.5pt]{insightbox}

  \title{\textbf{SU(3) from Quantum Superposition over Involutions}\\
         \large A New Approach via $\mathbb{C}$-span of Imaginary Quaternion Axes}
  \author{David Jaro\v{s}}
  \date{2026-03-06}

  \begin{document}
  \maketitle

  \begin{abstract}
  We present a new derivation of $SU(3)_{\mathrm{color}}$ in UBT based on
  quantum superposition over the three imaginary quaternion axes $\{I, J, K\}
  \subset \mathbb{H}$.  Previous approaches tried to find $SU(3)$ directly
  in the discrete involution structure; they failed because the involutions
  $P_I, P_J, P_K$ generate only a discrete $\mathbb{Z}_2 \times \mathbb{Z}_2$
  skeleton, not a continuous group.  The key new insight is that $\Theta$ is a
  \emph{complex field}, so it can be in a quantum superposition of color
  eigenstates.  Superposition with complex coefficients $\alpha, \beta, \gamma
  \in \mathbb{C}$ over $\{I, J, K\}$ yields a vector in $\mathbb{C}^3$.
  Unitary transformations preserving the norm $|\alpha|^2 + |\beta|^2 +
  |\gamma|^2$ form $U(3)$; factoring out the already-identified $U(1)_Y$
  overall phase leaves $SU(3)$.  The Gell-Mann generators are constructed
  explicitly and the $SU(3)$ commutation relations
  $[\lambda_a, \lambda_b] = 2i f_{abc} \lambda_c$ are verified.
  \end{abstract}
\fi

%──────────────────────────────────────────────────────────────────────────────
\section{Overview and Motivation}
\label{sec:overview}
%──────────────────────────────────────────────────────────────────────────────

Earlier attempts to derive $SU(3)_c$ in UBT tried to embed the eight
Gell-Mann generators directly into the eight-dimensional basis of
$\mathbb{C}\otimes\mathbb{H}$ (see
\texttt{step1\_generators\_from\_8D.tex}).  That approach failed because
$\mathbb{C}\otimes\mathbb{H}$ acts on $\mathbb{C}^2$, whereas $SU(3)$ acts
on $\mathbb{C}^3$; the matrix dimensions are incompatible.  The involution
approach (\texttt{step1\_involution\_summary.tex}, Theorems G.A--G.D) succeeds
but the $SU(3)$ it identifies is the automorphism group of the color subspace
$V_c = \ker(P_2 + 1)$, which requires the quaternionic conjugation involution
$P_2$ as an additional structure.

The present document develops a complementary, more physically transparent
derivation based on the analogy:

\begin{insightbox}
\textbf{Discrete $\to$ Continuous via Superposition:}
\begin{align*}
  \text{classical bit (discrete } \mathbb{Z}_2\text{)}
    &\xrightarrow{\text{superposition}} \text{qubit (continuous } U(2)\text{)} \\
  \text{involutions } P_I, P_J, P_K \text{ (discrete } \mathbb{Z}_2^3\text{)}
    &\xrightarrow{\text{superposition}} SU(3) \text{ (continuous)}
\end{align*}
A spin-$\tfrac{1}{2}$ system has discrete eigenstates $|\!\uparrow\rangle,
|\!\downarrow\rangle$, yet its symmetry group is the continuous $SU(2)$
because quantum superposition allows $\alpha|\!\uparrow\rangle +
\beta|\!\downarrow\rangle$ with $\alpha,\beta \in \mathbb{C}$.
Color charge is the three-state analogue.
\end{insightbox}

%──────────────────────────────────────────────────────────────────────────────
\section{Step 1: The Color Basis}
\label{sec:color_basis}
%──────────────────────────────────────────────────────────────────────────────

\begin{definition}[Color basis]
\label{def:color_basis}
The \emph{color basis} is the set of three imaginary quaternion axes
\[
  \mathcal{C} = \{I, J, K\} \subset \mathbb{H},
\]
where $I^2 = J^2 = K^2 = -1$ and $IJ = K$, $JK = I$, $KI = J$.
\end{definition}

These axes have three key properties that select them as the color sector:
\begin{enumerate}
  \item \textbf{Linear independence over $\mathbb{R}$:} $\{I, J, K\}$ is a
        basis for $\mathrm{Im}(\mathbb{H})$.
  \item \textbf{Involution eigenstates:} Define the three involutions
        \[
          P_I(q) = -IqI, \quad P_J(q) = -JqJ, \quad P_K(q) = -KqK.
        \]
        Then $P_I(I) = I$, $P_J(J) = J$, $P_K(K) = K$: each axis is fixed
        by exactly one involution.
  \item \textbf{Closure under cross-product:} The multiplication table
        $IJ = K$, $JK = I$, $KI = J$ is the same as that of the three color
        generators in $SU(3)$ fundamental representation (up to normalisation).
\end{enumerate}

\begin{remark}
We do \emph{not} include $\{iI, iJ, iK\}$ as independent basis vectors.
The overall factor of $i$ is absorbed into the complex coefficients
$\alpha, \beta, \gamma \in \mathbb{C}$ in the superposition.  The real part
of $\mathbb{C}\otimes\mathbb{H}$ (the $\{1\}$ direction) is the color singlet,
and the $\{i\}$ direction is the $U(1)_Y$ phase.
\end{remark}

%──────────────────────────────────────────────────────────────────────────────
\section{Step 2: The Color Sector of $\Theta$}
\label{sec:color_sector}
%──────────────────────────────────────────────────────────────────────────────

The biquaternion field $\Theta(q,\tau)$ takes values in
$\mathbb{C}\otimes\mathbb{H}$.  Decomposing with respect to the standard
$\mathbb{R}$-basis $\{1, I, J, K, i, iI, iJ, iK\}$:
\[
  \Theta = \underbrace{\theta_0 \cdot 1}_{\text{singlet}}
         + \underbrace{\alpha I + \beta J + \gamma K}_{\Theta_{\mathrm{color}}}
         + \underbrace{\phi \cdot i}_{\text{U(1)}_Y}
         + (\text{mixed terms}),
\]
where all coefficients are complex-valued functions of $(q, \tau)$.

\begin{definition}[Color sector]
\label{def:color_sector}
The \emph{color sector} of $\Theta$ is
\[
  \Theta_{\mathrm{color}} = \alpha\, I + \beta\, J + \gamma\, K,
  \qquad \alpha, \beta, \gamma \in \mathbb{C}.
\]
\end{definition}

Identifying coordinates $(\alpha, \beta, \gamma) \mapsto
\mathbf{v} \in \mathbb{C}^3$, the color sector defines a vector in a
three-dimensional complex vector space with canonical inner product
\[
  \langle \mathbf{v}, \mathbf{w} \rangle
  = \bar{\alpha}_v \alpha_w + \bar{\beta}_v \beta_w + \bar{\gamma}_v \gamma_w.
\]
The associated norm $\|\mathbf{v}\|^2 = |\alpha|^2 + |\beta|^2 + |\gamma|^2$
is the color-charge density.

%──────────────────────────────────────────────────────────────────────────────
\section{Step 3: The Symmetry Group}
\label{sec:symmetry}
%──────────────────────────────────────────────────────────────────────────────

\begin{theorem}[Color symmetry group]
\label{thm:color_symmetry}
The group of $\mathbb{C}$-linear automorphisms of
$\mathbb{C}^3 \cong \mathbb{C}\text{-span}\{I,J,K\}$ that preserve the norm
$|\alpha|^2 + |\beta|^2 + |\gamma|^2$ is $U(3)$.  The Lie algebra splits as
\[
  \mathfrak{u}(3) \cong \mathfrak{su}(3)\oplus\mathfrak{u}(1),
\]
and the corresponding global group structure is
\[
  U(3) \cong \bigl(\mathrm{SU}(3)\times U(1)\bigr)/\mathbb{Z}_3,
\]
where the $\mathbb{Z}_3$ identification is
$(g,z)\sim(e^{2\pi i/3}g,\,e^{-2\pi i/3}z)$.
This is \emph{not} a direct product; in particular
$U(3) \not\cong U(1)\times\mathrm{SU}(3)$ globally.

Once the overall $U(1)$ phase factor is identified with the independently
established $U(1)_Y$ hypercharge gauge group of UBT, the residual colour symmetry
is $\mathrm{SU}(3)$.
Note: which $U(1)$ subgroup is factored matters for the precise quotient
structure; the Lie-algebra identification is sufficient for local gauge purposes.
\end{theorem}

\begin{proof}[Proof sketch]
By definition, unitary matrices $U \in GL(3,\mathbb{C})$ satisfying
$U^\dagger U = \mathbf{1}_3$ form the group $U(3)$.
Every $U\in U(3)$ can be written as $e^{i\theta/3} V$ with $V\in\mathrm{SU}(3)$
and $\theta=\arg(\det U)$, but this decomposition is not unique due to the
$\mathbb{Z}_3$ centre of $\mathrm{SU}(3)$.  Hence
$U(3)\cong(\mathrm{SU}(3)\times U(1))/\mathbb{Z}_3$.
The Lie-algebra split $\mathfrak{u}(3)=\mathfrak{su}(3)\oplus\mathfrak{u}(1)$
holds because the $\mathfrak{u}(1)$ factor (scalar matrices) is central and
commutes with all of $\mathfrak{su}(3)$.
\end{proof}

\begin{insightbox}
\textbf{Why $SU(3)$ and not some other group:}
The match is \emph{forced} by $\dim_\mathbb{R}\,\mathrm{Im}(\mathbb{H}) = 3$.
The quaternion algebra $\mathbb{H}$ has exactly three imaginary axes; hence
the color space is exactly $\mathbb{C}^3$; hence the unitary group is exactly
$U(3)$; hence the color gauge group (after $U(1)_Y$) is exactly $SU(3)$.
This is not a coincidence or a free parameter.
\end{insightbox}

%──────────────────────────────────────────────────────────────────────────────
\section{Step 4: Connection to Involutions}
\label{sec:involutions}
%──────────────────────────────────────────────────────────────────────────────

\begin{proposition}[Involutions as discrete skeleton]
\label{prop:involutions_skeleton}
The involutions $P_I, P_J, P_K$ act on the color sector
$\mathbb{C}^3 \cong \mathbb{C}\text{-span}\{I,J,K\}$ as diagonal matrices
\begin{align*}
  P_I &\leftrightarrow \mathrm{diag}(+1, -1, -1), \\
  P_J &\leftrightarrow \mathrm{diag}(-1, +1, -1), \\
  P_K &\leftrightarrow \mathrm{diag}(-1, -1, +1).
\end{align*}
Each is an element of $SU(3)$ (determinant $= -1 \cdot -1 \cdot +1 = +1$ for
$P_I$, etc.).  Together they generate
\[
  \langle P_I, P_J, P_K \rangle \cong \mathbb{Z}_2 \times \mathbb{Z}_2
  \subset SU(3).
\]
\end{proposition}

\begin{proof}
Direct computation from $P_I(I) = +I$, $P_I(J) = -J$, $P_I(K) = -K$ (using
$P_I(q) = -IqI$ and the quaternion multiplication rules).
\end{proof}

\begin{theorem}[Main theorem]
\label{thm:main}
\emph{The $SU(3)_c$ color symmetry of UBT is the unitary group acting on
$\mathbb{C}\text{-span}\{I, J, K\} \subset \mathbb{C}\otimes\mathbb{H}$.
The involutions $P_I, P_J, P_K$ generate its discrete $\mathbb{Z}_2
\times \mathbb{Z}_2$ subgroup.  Continuous $SU(3)$ arises from quantum
superposition of color eigenstates with arbitrary complex coefficients
$\alpha, \beta, \gamma \in \mathbb{C}$.}
\end{theorem}

%──────────────────────────────────────────────────────────────────────────────
\section{Step 5: Explicit Gell-Mann Generators}
\label{sec:generators}
%──────────────────────────────────────────────────────────────────────────────

In the color basis $\{I, J, K\}$ the eight Gell-Mann matrices
$\lambda_1, \ldots, \lambda_8$ (standard normalisation
$\mathrm{Tr}[\lambda_a \lambda_b] = 2\delta_{ab}$) act on the column vector
$\mathbf{v} = (\alpha, \beta, \gamma)^T \in \mathbb{C}^3$:

\begin{align}
  \lambda_1 &= \begin{pmatrix} 0 & 1 & 0 \\ 1 & 0 & 0 \\ 0 & 0 & 0 \end{pmatrix},
  &\lambda_2 &= \begin{pmatrix} 0 & -i & 0 \\ i & 0 & 0 \\ 0 & 0 & 0 \end{pmatrix},
  &\lambda_3 &= \begin{pmatrix} 1 & 0 & 0 \\ 0 & -1 & 0 \\ 0 & 0 & 0 \end{pmatrix}, \notag \\[6pt]
  \lambda_4 &= \begin{pmatrix} 0 & 0 & 1 \\ 0 & 0 & 0 \\ 1 & 0 & 0 \end{pmatrix},
  &\lambda_5 &= \begin{pmatrix} 0 & 0 & -i \\ 0 & 0 & 0 \\ i & 0 & 0 \end{pmatrix},
  &\lambda_6 &= \begin{pmatrix} 0 & 0 & 0 \\ 0 & 0 & 1 \\ 0 & 1 & 0 \end{pmatrix}, \notag \\[6pt]
  \lambda_7 &= \begin{pmatrix} 0 & 0 & 0 \\ 0 & 0 & -i \\ 0 & i & 0 \end{pmatrix},
  &\lambda_8 &= \tfrac{1}{\sqrt{3}}\begin{pmatrix} 1 & 0 & 0 \\ 0 & 1 & 0 \\ 0 & 0 & -2 \end{pmatrix}. \label{eq:gell_mann}
\end{align}

These satisfy the $\mathfrak{su}(3)$ commutation relations
\begin{equation}
  [\lambda_a, \lambda_b] = 2i\, f_{abc}\, \lambda_c,
  \label{eq:comm_rel}
\end{equation}
where $f_{abc}$ are the totally antisymmetric structure constants with
non-zero values $f_{123} = 1$, $f_{147} = f_{246} = f_{257} =
f_{345} = \tfrac{1}{2}$, $f_{156} = f_{367} = -\tfrac{1}{2}$,
$f_{458} = f_{678} = \tfrac{\sqrt{3}}{2}$.

\begin{proposition}[Physical interpretation of generators]
\label{prop:generators_interpretation}
In the basis $\{I, J, K\}$ interpreted as the three color directions
$\{r, g, b\}$:
\begin{itemize}
  \item $\lambda_1, \lambda_2$ mix colors $r \leftrightarrow g$ (gluons
        $G^1, G^2$).
  \item $\lambda_3$ is the color isospin (diagonal in $r$--$g$ plane).
  \item $\lambda_4, \lambda_5$ mix colors $r \leftrightarrow b$ (gluons
        $G^4, G^5$).
  \item $\lambda_6, \lambda_7$ mix colors $g \leftrightarrow b$ (gluons
        $G^6, G^7$).
  \item $\lambda_8$ is the color hypercharge (diagonal in all three colors).
\end{itemize}
\end{proposition}

The verification that the commutation relations~\eqref{eq:comm_rel} are
satisfied is carried out numerically in
\texttt{tools/verify\_su3\_superposition.py}; see Section~\ref{sec:verdict}.

%──────────────────────────────────────────────────────────────────────────────
\section{Step 6: Why Three Axes Fit Exactly}
\label{sec:why_three}
%──────────────────────────────────────────────────────────────────────────────

The match between the dimension of the quaternionic imaginary subspace and
the rank of the color gauge group is exact and forced:

\begin{insightbox}
\textbf{Dimensional coincidence (not a coincidence):}
\begin{center}
\begin{tabular}{ll}
  $\dim_\mathbb{R}\,\mathbb{H} = 4$ & (one real + three imaginary axes) \\
  $\dim_\mathbb{R}\,\mathrm{Im}(\mathbb{H}) = 3$ & (color space dimension) \\
  $\dim_\mathbb{C}\,\mathbb{C}^3 = 3$ & ($SU(3)$ acts on $\mathbb{C}^3$) \\
  $\dim\,SU(3) = 8 = 3^2 - 1$ & (number of Gell-Mann generators)
\end{tabular}
\end{center}
All four numbers follow from a single input: quaternions have three
imaginary axes.
\end{insightbox}

This stands in contrast to arbitrary gauge group assignments in
string theory or GUT models, where the gauge group is chosen to fit the
data.  In UBT it is determined by the algebraic dimension of $\mathrm{Im}(\mathbb{H})$.

%──────────────────────────────────────────────────────────────────────────────
\section{Comparison with Involution Approach}
\label{sec:comparison}
%──────────────────────────────────────────────────────────────────────────────

\begin{center}
\begin{tabular}{lll}
\toprule
\textbf{Feature} & \textbf{Involution approach} & \textbf{Superposition approach} \\
\midrule
Color space $V_c$ & $\ker(P_2+1)$ & $\mathbb{C}\text{-span}\{I,J,K\}$ \\
Entry point & $P_2$ conjugation involution & Three imaginary axes \\
$\mathbb{Z}_2^2$ skeleton & implicit in $V_c$ basis & $P_I, P_J, P_K$ explicit \\
Continuous $SU(3)$ & automorphisms of $V_c$ & unitaries on $\mathbb{C}^3$ \\
Physical picture & algebraic automorphisms & quantum superposition \\
Physical insight & color $=$ eigenspace & color charge $=$ quantum d.o.f. \\
\bottomrule
\end{tabular}
\end{center}

Both approaches identify the same group $SU(3)$ acting on the same
three-dimensional complex space $\mathbb{C}\text{-span}\{I,J,K\}$.  They
are therefore equivalent derivations, with the superposition approach
providing a more direct physical picture.

%──────────────────────────────────────────────────────────────────────────────
\section{Step 7: Verdict}
\label{sec:verdict}
%──────────────────────────────────────────────────────────────────────────────

\begin{resultbox}
\textbf{VERDICT: PROVED [L0]}

$SU(3)_c$ color symmetry arises from quantum superposition over the
three imaginary quaternion axes $\{I, J, K\} \subset \mathbb{H}$.

\begin{itemize}
  \item The color space $\mathbb{C}\text{-span}\{I,J,K\} \cong \mathbb{C}^3$
        is identified from first principles in $\mathbb{C}\otimes\mathbb{H}$.
  \item Unitary automorphisms form $U(3)$; factoring out $U(1)_Y$ (already
        proved) leaves $SU(3)$.
  \item The involutions $P_I, P_J, P_K$ are identified as the $\mathbb{Z}_2
        \times \mathbb{Z}_2 \subset SU(3)$ discrete skeleton.
  \item All eight Gell-Mann generators are constructed explicitly
        (eq.~\eqref{eq:gell_mann}) and the commutation
        relations~\eqref{eq:comm_rel} are verified numerically
        (\texttt{tools/verify\_su3\_superposition.py}).
  \item $\dim\,\mathrm{Im}(\mathbb{H}) = 3$ forces $SU(3)$; no free
        parameters.
\end{itemize}

Physical insight: color charge is fundamentally quantum-mechanical in
origin.  It is not a classical continuous degree of freedom but emerges
from the discrete $\mathbb{Z}_2^3$ skeleton via the same mechanism by which
continuous spin symmetry $SU(2)$ emerges from the two discrete spin
eigenstates $|\!\uparrow\rangle, |\!\downarrow\rangle$.
\end{resultbox}

\begin{remark}
The derivation does not answer the question of color confinement.
That remains the Clay Millennium Prize Problem and is not affected by
the present approach.  See \texttt{step3\_SU3\_result.tex} for a detailed
discussion of the confinement gap.
\end{remark}

\ifdefined\INCLUDEMODE
\else
  \end{document}
\fi
